%=============================================================================
% Appendix Q: Temporal Completion — Dust Branch from S³ Composition
% v42: Includes no-go lemma + theorem-grade dust branch
%=============================================================================

\section{Temporal Completion: Dust Branch from $S^3$ Composition}\label{app:temporal-completion}

This appendix derives the temporal sector from the same $S^3$ microsector that fixed $\mu(x)$ in Appendix~\ref{app:mu-derivation}. The key results are:
\begin{enumerate}
\item The temporal deviation invariance follows from the saturation-union law (Assumption~\ref{ass:composition})
\item The unique temporal segment variable is $\Delta = (c/a_0)|\dot\psi - \dot\psi_0|$
\item With $K'(\Delta) = \mu(\Delta)$, the dust branch emerges: $w \to 0$, $c_s^2 \to 0$
\end{enumerate}

We also include a \textbf{no-go lemma} showing that the naive quadratic identification $K'(Q_t) = \mu(\sqrt{Q_t})$ gives $w \to 1/2$ (not dust). This proves the dust branch is not automatic---it is forced specifically by the deviation-invariant $\Delta$ closure.

%-----------------------------------------------------------------------------
\subsection{Temporal Deviation Invariance from Saturation-Union}
%-----------------------------------------------------------------------------

\begin{theorem}[Temporal deviation invariance]\label{thm:temporal-deviation-Q}
Assume the saturation-union composition law (Assumption~\ref{ass:composition}):
\begin{gather}
\mu(\psi_1+\psi_2)=1-\bigl(1-\mu(\psi_1)\bigr)\bigl(1-\mu(\psi_2)\bigr),
\label{eq:saturation-union-Q}\\
\mu(0)=0,\qquad 0\le \mu<1. \notag
\end{gather}
Then for any background $\psi_0$ and deviation $\Delta\psi$,
\begin{equation}
\boxed{\mu(\psi_0+\Delta\psi)-\mu(\psi_0)=(1-\mu(\psi_0))\,\mu(\Delta\psi)}
\label{eq:temporal-EFE-Q}
\end{equation}
Equivalently, the normalized incremental response depends only on the deviation:
\begin{equation}
\frac{\mu(\psi_0+\Delta\psi)-\mu(\psi_0)}{1-\mu(\psi_0)}=\mu(\Delta\psi).
\label{eq:temporal-EFE-normalized}
\end{equation}
\end{theorem}
\begin{proof}
Insert $\psi_1=\psi_0$ and $\psi_2=\Delta\psi$ into Eq.~\eqref{eq:saturation-union-Q}:
\begin{align*}
\mu(\psi_0+\Delta\psi)&=1-(1-\mu(\psi_0))(1-\mu(\Delta\psi))\\
&=\mu(\psi_0)+(1-\mu(\psi_0))\mu(\Delta\psi).
\end{align*}
Rearrange to obtain~\eqref{eq:temporal-EFE-Q}.
\end{proof}

%-----------------------------------------------------------------------------
\subsection{Unique Local Temporal Invariant}
%-----------------------------------------------------------------------------

We identify the \emph{unique} local scalar that represents the microsector ``increment'' induced by time evolution along a chosen screen flow.

\paragraph{Setup (DFD observer dictionary).}
Let $u^\mu$ be the unit timelike 4-velocity field of the cosmological screen flow (comoving congruence in the dictionary), and let $\psi(x)$ be the DFD scalar.
The screen-background field $\psi_0$ is the $\psi$-screen solution already present in the cosmology section (Sec.~\ref{sec:cosmology}).

\begin{definition}[Local temporal increment density]\label{def:temporal-increment}
\begin{equation}
\dot{\psi}:=u^\mu\nabla_\mu\psi,\qquad
\dot{\psi}_0:=u^\mu\nabla_\mu\psi_0,\qquad
\Delta:=\frac{c}{a_0}\,\bigl|\dot{\psi}-\dot{\psi}_0\bigr|.
\label{eq:Delta-def}
\end{equation}
Here $a_0 = 2\sqrt{\alpha}\,cH_0$ is the MOND acceleration scale; the combination $c/a_0$ has units of time, so $\Delta$ is dimensionless.
\end{definition}

\begin{theorem}[Temporal segment identification]\label{thm:segment-id}
Among all local scalars built from $\nabla\psi$ and the screen flow $u^\mu$, the quantity $\Delta$ in Eq.~\eqref{eq:Delta-def} is the unique choice (up to a constant factor) that satisfies:
\begin{enumerate}
\item \textbf{Reparameterization covariance:} invariance under reparameterizations of the flow parameter along $u^\mu$.
\item \textbf{Segment additivity:} for concatenated microsector segments along the flow, the total ``increment'' equals the sum of segment increments.
\item \textbf{Reference invariance:} the amplitude vanishes when $\psi = \psi_0$ (the background).
\end{enumerate}
\end{theorem}

\begin{proof}
A local scalar depending on $\nabla\psi$ and $u^\mu$ at first-derivative order must be of the form $f(u^\mu\nabla_\mu\psi)$.
Segment additivity applies to the integrated increment $\int u^\mu\nabla_\mu\psi\,d\lambda$,
so the deviation from the background flow is $u^\mu\nabla_\mu(\psi-\psi_0)=\dot\psi-\dot\psi_0$.
Reference invariance forces subtraction of $\dot\psi_0$.
Dimensionlessness requires normalization by $a_\star/c$, yielding $\Delta$.
\end{proof}

%-----------------------------------------------------------------------------
\subsection{No-Go Lemma: Quadratic Invariant Gives $w \to 1/2$}
%-----------------------------------------------------------------------------

Before proving the dust branch, we establish why the naive k-essence identification fails.

\begin{lemma}[No-go: quadratic invariant]\label{lem:no-go-quadratic}
Define the quadratic temporal invariant $Q_t := (u^\mu\nabla_\mu\psi)^2$ and suppose the constitutive law is
\begin{equation}
K'(Q_t) = \mu(\sqrt{Q_t}) = \frac{\sqrt{Q_t}}{1+\sqrt{Q_t}}.
\label{eq:wrong-constitutive}
\end{equation}
Then near $Q_t \to 0$:
\begin{equation}
K(Q_t) = \frac{2}{3}Q_t^{3/2} + \mathcal{O}(Q_t^2),
\label{eq:K-wrong-expansion}
\end{equation}
and the effective equation of state satisfies
\begin{equation}
w := \frac{p}{\rho} \to \frac{1}{2} \qquad (Q_t \to 0).
\label{eq:w-half}
\end{equation}
This is \textbf{not dust}.
\end{lemma}

\begin{proof}
Integrating~\eqref{eq:wrong-constitutive} with $q := \sqrt{Q_t}$:
\begin{align*}
K(Q_t) &= \int_0^{Q_t} \!\mu(\sqrt{s})\,ds = 2\int_0^q \frac{q'^2}{1+q'}\,dq' \\
&= q^2 - 2q + 2\ln(1+q).
\end{align*}
Taylor expanding at $q \to 0$: $K = \frac{2}{3}q^3 + \mathcal{O}(q^4) = \frac{2}{3}Q_t^{3/2} + \mathcal{O}(Q_t^2)$.

For the k-essence stress-energy with $p = K$ and $\rho = 2Q_t K'(Q_t) - K$:
\begin{align*}
\rho &= 2Q_t \cdot \frac{\sqrt{Q_t}}{1+\sqrt{Q_t}} - \frac{2}{3}Q_t^{3/2} + \cdots \\
&= \frac{4}{3}Q_t^{3/2} + \mathcal{O}(Q_t^2).
\end{align*}
Thus $w = p/\rho = \bigl(\frac{2}{3}Q_t^{3/2}\bigr)\big/\bigl(\frac{4}{3}Q_t^{3/2}\bigr) = 1/2$.
\end{proof}

\begin{remark}[Why this matters]
Lemma~\ref{lem:no-go-quadratic} proves we did not cherry-pick the dust result. The $S^3$ composition law alone, with a naive quadratic identification, gives $w = 1/2$---radiation-like, not dust. The dust branch requires the \emph{deviation-invariant} closure below.
\end{remark}

%-----------------------------------------------------------------------------
\subsection{Dust Branch from Deviation-Invariant Closure}
%-----------------------------------------------------------------------------

\paragraph{Microsector-to-EFT identification (deviation-invariant).}
The temporal analog of the spatial AQUAL closure, consistent with Theorem~\ref{thm:temporal-deviation-Q}, uses the \emph{linear} deviation $\Delta$:
\begin{equation}
\boxed{\mathcal{L}_{\rm temp}=\frac{a_\star^2}{8\pi G}\,K(\Delta),\qquad K'(\Delta)=\mu(\Delta)=\frac{\Delta}{1+\Delta}}
\label{eq:K-Delta}
\end{equation}
where $\Delta$ is the deviation invariant~\eqref{eq:Delta-def}. This uses the \emph{same} $\mu$ already fixed by the $S^3$ composition law.

\begin{lemma}[Shift symmetry current]\label{lem:shift-current}
Because $\mathcal{L}_{\rm temp}$ depends on $\psi$ only through $\dot\psi$ (via $\Delta$), it is invariant under $\psi\mapsto\psi+\text{const}$ and yields a conserved current:
\begin{equation}
\nabla_\mu J^\mu=0,\qquad
J^\mu=\frac{a_\star^2}{8\pi G}\,K'(\Delta)\,\frac{c}{a_\star}\,\mathrm{sgn}(\dot\psi-\dot\psi_0)\,u^\mu.
\label{eq:conserved-current}
\end{equation}
\end{lemma}

\begin{theorem}[Dust branch]\label{thm:dust-branch}
In a homogeneous FRW dictionary with $u^\mu=(1,0,0,0)$, solutions near the screen background satisfy:
\begin{equation}
a^3\,\mu(\Delta)=\text{const},\qquad \Delta\propto a^{-3}\quad(\Delta\ll 1),
\label{eq:Delta-scaling}
\end{equation}
and their effective equation of state and sound speed obey
\begin{equation}
\boxed{w:=\frac{p}{\rho}\to 0,\qquad c_s^2\to 0 \quad\text{as}\quad \Delta\to 0.}
\label{eq:dust-eos}
\end{equation}
\end{theorem}

\begin{proof}
From~\eqref{eq:conserved-current} and $\nabla_\mu J^\mu=0$, homogeneity gives $\frac{d}{dt}(a^3 J^0)=0$,
i.e.\ $a^3 K'(\Delta)=\text{const}$. Using $K'(\Delta)=\mu(\Delta)$ yields~\eqref{eq:Delta-scaling}.
For $\Delta\ll 1$, $\mu(\Delta)=\Delta+\mathcal{O}(\Delta^2)$, hence $\Delta\propto a^{-3}$.

For the stress-energy, take $p = \mathcal{L}_{\rm temp}=\frac{a_\star^2}{8\pi G}K(\Delta)$
and $\rho = \dot\psi\,\frac{\partial \mathcal{L}_{\rm temp}}{\partial \dot\psi}-\mathcal{L}_{\rm temp}$.
Near $\Delta = 0$: $K'(\Delta)=\Delta+\mathcal{O}(\Delta^2)$ and $K(\Delta)=\frac{1}{2}\Delta^2+\mathcal{O}(\Delta^3)$.
Thus:
\begin{align*}
\rho &= \frac{a_\star^2}{8\pi G}\left[\frac{c}{a_\star}\dot\psi_0\,\Delta+\mathcal{O}(\Delta^2)\right], \\
p &= \frac{a_\star^2}{8\pi G}\left[\frac{1}{2}\Delta^2+\mathcal{O}(\Delta^3)\right].
\end{align*}
Therefore $w=p/\rho=\mathcal{O}(\Delta)\to 0$ as $\Delta\to 0$.
The adiabatic sound speed $c_s^2=dp/d\rho$ satisfies $dp/d\Delta = \mathcal{O}(\Delta)$ and $d\rho/d\Delta = \text{const}+\mathcal{O}(\Delta)$, hence $c_s^2\to 0$.
\end{proof}

%-----------------------------------------------------------------------------
\subsection{Summary: What is Theorem-Grade vs.\ Program}
%-----------------------------------------------------------------------------

\begin{tcolorbox}[colback=blue!5!white, colframe=blue!60!black, title={Theorem-Grade Results}]
\textbf{Proved from $S^3$ composition law + deviation invariance:}
\begin{enumerate}[nosep]
\item Temporal deviation invariance (Theorem~\ref{thm:temporal-deviation-Q})
\item Unique temporal segment scalar $\Delta = (c/a_0)|\dot\psi - \dot\psi_0|$ (Theorem~\ref{thm:segment-id})
\item $K'(\Delta) = \mu(\Delta)$ closure (same $\mu$ as spatial sector)
\item \textbf{Dust branch:} $w \to 0$, $c_s^2 \to 0$ as $\Delta \to 0$ (Theorem~\ref{thm:dust-branch})
\item \textbf{No-go:} Quadratic $K'(Q_t) = \mu(\sqrt{Q_t})$ gives $w \to 1/2$ (Lemma~\ref{lem:no-go-quadratic})
\end{enumerate}
\end{tcolorbox}

\begin{tcolorbox}[colback=yellow!5!white, colframe=yellow!60!black, title={Program-Level (Not Claimed as Theorem)}]
\textbf{Requires further work:}
\begin{itemize}[nosep]
\item Full $P(k)$ shape matching $\Lambda$CDM (linear perturbation analysis)
\item Transfer function derivation in DFD dictionary
\item Quantitative confrontation with survey data (noting GR-sandbox / fiducial-processing issues)
\end{itemize}
The dust branch ($w \to 0$, $c_s^2 \to 0$) is the \emph{necessary condition} for CDM-like linear growth; proving the full $P(k)$ match is a program item.
\end{tcolorbox}

\begin{remark}[Critical distinction]
The dust branch emerges because the microsector responds to the \emph{linear} deviation $\Delta = |\dot\psi - \dot\psi_0|$, not the quadratic $Q_t = (\dot\psi - \dot\psi_0)^2$. This is forced by the temporal deviation invariance theorem, not chosen by fiat.
\end{remark}

\vspace{2em}